% Part: first-order-logic
% Chapter: sequent-calculus
% Section: provability-consistency

\documentclass[../../../include/open-logic-section]{subfiles}

\begin{document}

\iftag{FOL}
      {\olfileid{fol}{seq}{prv}}
      {\olfileid{pl}{seq}{prv}}

\olsection{\usetoken{S}{derivability} lan Konsistensi}

Saiki kita bakal netepake sawetara sipat relasi !!{derivability}.
Sipat kasebut narik kawigaten dhewe-dhewe, nanging saben sipat bakal
duwe peran ing bukti teorema kelengkapan.

\begin{prop}\ollabel{prop:provability-contr}
  Yen $\Gamma \Proves !A$ lan $\Gamma \cup \{!A\}$ ora konsisten,
  mula $\Gamma$ ora konsisten.
\end{prop}

\begin{proof}
Ana $\Gamma_0$ lan $\Gamma_1$ kang mènèhi !!{derivation}s
$\pi_0$ kanggo $\Gamma_0 \Sequent !A$ lan $\pi_1$ kanggo
$!A, \Gamma_1 \Sequent \quad$. Subhimpunan loro mau winates lan kalebu ing $\Gamma$.
Pelemahan lan pertukaran kiwa bisa digunakaké kanggo mapanaké
$!A$ ing pucuk kiwa premis kapindho. Banjur cut mènèhi
\begin{prooftree}
\AxiomC{}
\RightLabel{$\pi_0$}
\Deduce$ \Gamma_0 \fCenter !A $
\AxiomC{}
\RightLabel{$\pi_1$}
\Deduce$!A, \Gamma_1 \fCenter $
\RightLabel{\Cut}
\BinaryInf$ \Gamma_0,\Gamma_1 \fCenter $
\end{prooftree}

Amarga $\Gamma_0 \subseteq \Gamma$ lan $\Gamma_1 \subseteq \Gamma$,
$\Gamma_0 \cup \Gamma_1 \subseteq \Gamma$, mula $\Gamma$~ora
konsisten.
\end{proof}

\begin{prop}
\ollabel{prop:prov-incons}
$\Gamma \Proves !A$ yen lan mung yen $\Gamma \cup \{\lnot !A\}$ ora
konsisten.
\end{prop}

\begin{proof}
Sepisan, upamakna $\Gamma \Proves !A$. Ana !!a{derivation}
$\pi_0$ kanggo $\Gamma_0 \Sequent !A$, kanthi $\Gamma_0$ subhimpunan
winates saka $\Gamma$. Kanthi aturan $\LeftR{\lnot}$, kita entuk
!!a{derivation} saka $\lnot !A, \Gamma_0 \Sequent \quad$.
Mula $\Gamma \cup \{\lnot !A\}$ ora konsisten.

Yèn $\Gamma \cup \{\lnot !A\}$ ora konsisten, jupuken subhimpunan
winates $\Gamma_1$ saka $\Gamma$ lan !!a{derivation} $\pi_1$ saka
$\lnot !A, \Gamma_1 \Sequent \quad$ banjur sawisé pelemahan lan pertukaran kiwa
yèn perlu. Wit iki mènèhi $\Gamma_1 \Sequent !A$.
\begin{prooftree}
  \Axiom$!A \fCenter !A$
  \RightLabel{\RightR{\lnot}}
  \UnaryInf$\fCenter !A, \lnot !A$
  \AxiomC{}
  \RightLabel{$\pi_1$}
  \Deduce$\lnot !A, \Gamma_1 \fCenter$
  \RightLabel{\Cut}
  \BinaryInf$\Gamma_1 \fCenter !A$
\end{prooftree}
Amarga $\Gamma_1 \subseteq \Gamma$, kita éntuk $\Gamma \Proves !A$.
\end{proof}

\begin{prob}
Buktekna yen $\Gamma \Proves \lnot !A$ yen lan mung yen $\Gamma \cup
\{!A\}$ ora konsisten.
\end{prob}

\begin{prop}\ollabel{prop:explicit-inc}
  Yen $\Gamma \Proves !A$ lan $\lnot !A \in \Gamma$, mula $\Gamma$
  ora konsisten.
\end{prop}

\begin{proof}
  Upamana $\Gamma \Proves !A$ lan $\lnot !A \in \Gamma$. Banjur ana
  !!a{derivation}~$\pi$ saka sekuen $\Gamma_0 \Sequent !A$, kanthi
  $\Gamma_0$ subhimpunan winates saka $\Gamma$. Sekuen
  $\Gamma_0, \lnot !A \Sequent \quad$ uga !!{derivable}:
  \begin{prooftree}
    \AxiomC{}
    \RightLabel{$\pi$}
    \Deduce$\Gamma_0 \fCenter !A$
    \Axiom$!A \fCenter !A$
    \RightLabel{\LeftR{\lnot}}
    \UnaryInf$\lnot !A, !A \fCenter$
    \RightLabel{\LeftR{\Exchange}}
    \UnaryInf$!A, \lnot !A \fCenter$
    \RightLabel{\Cut}
    \BinaryInf$\Gamma_0, \lnot !A \fCenter$
  \end{prooftree}
  Amarga $\lnot !A \in \Gamma$ lan $\Gamma_0 \subseteq \Gamma$, iki
  nuduhake yen $\Gamma$ ora konsisten.
\end{proof}

\begin{prop}\ollabel{prop:provability-exhaustive}
  Yen $\Gamma \cup \{!A\}$ lan $\Gamma \cup \{\lnot !A\}$ kalorone
  ora konsisten, mula $\Gamma$ ora konsisten.
\end{prop}

\begin{proof}
Ana himpunan winates $\Gamma_0 \subseteq \Gamma$ lan $\Gamma_1
\subseteq \Gamma$, sarta $\Log{LK}$-!!{derivation}s $\pi_0$ lan
$\pi_1$ saka $!A, \Gamma_0 \Sequent \quad$ lan $\lnot !A, \Gamma_1
\Sequent \quad$, kanthi urut. Pelemahan lan pertukaran kiwa njamin
panggonané formula kapisan ing saben premis. Banjur kita bisa !!{derive}
\begin{prooftree}
\AxiomC{}
\RightLabel{$\pi_0$}
\Deduce$ !A, \Gamma_0 \fCenter $
\RightLabel{\RightR{\lnot}}
\UnaryInf$ \Gamma_0 \fCenter \lnot !A$
\AxiomC{}
\RightLabel{$\pi_1$}
\Deduce$\lnot !A, \Gamma_1 \fCenter  $
\RightLabel{\Cut}
\BinaryInf$ \Gamma_0, \Gamma_1 \fCenter $
\end{prooftree}
Amarga $\Gamma_0 \subseteq \Gamma$ lan $\Gamma_1 \subseteq \Gamma$,
$\Gamma_0 \cup \Gamma_1 \subseteq \Gamma$. Mula $\Gamma$ ora
konsisten.
\end{proof}

\end{document}
